% Original-source calculation. No auxiliary Gamma order is introduced.
\section{The intrinsic exterior mass and the original four-endpoint relation graph}
\label{sec:IFB}

The source order in this section is $s\in\{1,k\}$, with
$k=4l+1$, $l,m\geq1$, $e=1+k(m-1)$,
$q=e(k+1)^2$, and $c=k/2$.
Retain $0<\delta<1/2$, $\gamma>2$ and the complete original polynomial
\[
 \chi(S)=\prod_{a,b=0}^k
 [S-c-(2a-k)\delta-i(2b-k)\gamma]^e.
 \tag{IFB1}
\]
The repeated index $b$ in this product is local to the product; below
$b_s=s/2$ denotes the source parameter. We write
\[
 M_s=(2\pi)^{s/2},\qquad b_s=s/2,\qquad
 dm_{0,s}(y)=
 \frac{M_s2^{b_s-2}}{\pi\Gamma(b_s)}
 \left|\Gamma\left(\frac{b_s}{2}+\frac{iy}{2}\right)\right|^2dy.
 \tag{IFB2}
\]
For $s=1$ this is exactly $d\sigma$ in BSL1. For $s=k$ it is
exactly the original tensor-Gamma measure in BSL17: repeated use of
$\Gamma(z+1)=z\Gamma(z)$ gives
\[
 \left|\Gamma\left(l+\frac14+\frac{iy}{2}\right)\right|^2
 =4^{-l}\prod_{j=0}^{l-1}\bigl(y^2+(2j+\tfrac12)^2\bigr)
       |\Gamma(\tfrac14+\tfrac{iy}{2})|^2,
\]
and the coefficient in IFB2 times $4^{-l}$ equals
$M_k/(2\pi\sqrt2\Gamma(k/2))$. Thus no source has been selected
by matching its size to $q$. Both sources are the measures already
present in the original calculation. The original evaluation line in
both cases is $S=c+iy$.

\subsection{The original orthogonal coordinates and the complete quotient map}

The original source transform, with its full mass, is
\[
 \int_{\mathbb R}e^{ty}\,dm_{0,s}(y)
       =M_s(\cos t)^{-b_s},\qquad |t|<\pi/2.
 \tag{IFB3}
\]
Here is a direct check of its scale. With $\lambda=b_s/2$ and
$g(x)=\exp(\lambda x-e^x)$, Euler's integral gives
$\widehat g(\xi)=\Gamma(\lambda+i\xi)$ for
$\widehat g(\xi)=\int g(x)e^{i\xi x}dx$. Autocorrelation gives
\[
 \int e^{iu\xi}|\Gamma(\lambda+i\xi)|^2d\xi
 =2\pi\int g(x+u/2)g(x-u/2)dx
 =\frac{2\pi\Gamma(2\lambda)}{(2\cosh(u/2))^{2\lambda}}.
\]
The first equality follows by Fourier inversion of the integrable
autocorrelation. The function $g$ and every derivative decay
exponentially at $-\infty$ and faster than exponentially at $+\infty$,
so these Fourier operations and the displayed integrals converge.
The last equality follows from $v=e^x$. Substituting $y=2\xi$,
$u=2w$, and the coefficient in IFB2 gives
$M_s(\cosh w)^{-b_s}$. Rotation of Euler's integral through
$\theta\operatorname{sgn}(y)$, $0<\theta<\pi/2$, proves an
upper bound $C_{s,\theta}\exp(-\theta|y|)$ for the squared Gamma
factor. The small sector arc tends to zero because $\lambda>0$;
the large arc tends to zero because $\cos\theta>0$. This proves
analytic continuation of the characteristic transform to
$|\operatorname{Im}w|<\pi/2$ and hence IFB3. In particular all
polynomial moments below exist and $\int dm_{0,s}=M_s$.

Define real monic polynomials $p_n^{(s)}(y)$ by
\[
 \sum_{n=0}^\infty p_n^{(s)}(y)\frac{z^n}{n!}
 =(1+z^2)^{-b_s/2}\exp(y\arctan z).
 \tag{IFB4}
\]
Branches are analytic at zero. For sufficiently small real $z,w$,
IFB3 and the cosine addition formula show that the integral of the
product of these two generating functions is
$M_s(1-zw)^{-b_s}$. The exponential bound used for IFB3 permits
each derivative in a neighborhood of zero under the integral.
Coefficient comparison gives the exact norms and recurrence
\[
 h_n^{(s)}=M_s n!(b_s)_n,\qquad
 \int p_n^{(s)}p_j^{(s)}\,dm_{0,s}
       =\mathbf1_{n=j}h_n^{(s)},
\]
\[
 yp_n^{(s)}=p_{n+1}^{(s)}+n(b_s+n-1)p_{n-1}^{(s)},\quad
 p_0^{(s)}=1,\quad p_{-1}^{(s)}=0.
 \tag{IFB5}
\]
The recurrence also follows by differentiating IFB4 and multiplying
by $1+z^2$. Each $p_n^{(s)}$ has parity $n$. We henceforth retain
the superscript $s$ on $h_n$ where needed and suppress it on $p_n$.

Let $E_\chi=\mathbb C[S]/(\chi)$ with its fixed remainder frame
$(1,S,\ldots,S^{q-1})$, and define actual source polynomials
\[
 u_n(S)=\frac{p_n((S-c)/i)}{\sqrt{h_n^{(s)}}},\qquad
 C_s:\mathbb C^q\longrightarrow E_\chi,\quad
 C_s x=\sum_{n=0}^{q-1}x_n[u_n].
 \tag{IFB6}
\]
The positive square root is used. On the original line $u_n(c+iy)$
equals $p_n(y)/\sqrt{h_n^{(s)}}$. The source polynomials $u_n$ are
therefore an orthonormal frame, with all source factors explicitly
present. The leading coefficient of $u_n$ in $S$ is
$i^{-n}/\sqrt{h_n^{(s)}}$. Thus $C_s$ is triangular and invertible,
and
\[
 |\det C_s|^2=\prod_{n=0}^{q-1}(h_n^{(s)})^{-1}.
 \tag{IFB7}
\]
This is an exact change of coordinates on the original remainder
space, not a replacement of that space or of its frame.

For $j\geq0$ put
\[
 v_j=C_s^{-1}[u_{q+j}]\in\mathbb C^q,\qquad
 V_r=[v_0\ \cdots\ v_{r-1}],\qquad
 K_{s,r}=I_q+V_rV_r^*,\qquad \Omega_{s,r}=\det K_{s,r}.
 \tag{IFB8}
\]
The definitions for $r=0$ are $V_0$ the empty $q$ by zero matrix,
$K_{s,0}=I_q$, and $\Omega_{s,0}=1$. These vectors use the same
$C_s$ and the same remainder map for every endpoint. Fix
$r\geq0$ and $N=q+r-1$. The orthonormal source-coordinate map
\[
 U_{s,N}:\mathcal P_N\longrightarrow\mathbb C^q\oplus\mathbb C^r,
 \quad
 \sum_{n<q}x_nu_n+\sum_{j<r}z_ju_{q+j}\longmapsto(x,z)
\]
is a linear isometry. Polynomial division and IFB8 give the exact
commuting relation
\[
 J_NU_{s,N}^{-1}(x,z)=C_s(x+V_rz),\qquad
 U_{s,N}(\chi\mathcal P_{r-1})
       =\{(-V_rz,z):z\in\mathbb C^r\}.
 \tag{IFB9}
\]
To verify equality on the right, a vector in the displayed graph
has zero remainder by the first formula and hence is divisible by
the original monic $\chi$. Its degree is at most $q+r-1$, so its
quotient has degree at most $r-1$. Conversely every such multiple
has zero remainder, which forces its low coordinate to be $-V_rz$.
The map $z\mapsto U_{s,N}^{-1}(-V_rz,z)$ is consequently an exact
bijection from the high source coordinates to the original relation
space. No relation direction is projected out. When $r=0$, this
means the zero space on both sides.

For a fixed original remainder $a\in E_\chi$, write $w=C_s^{-1}a$.
The unique minimum source-norm lift has coordinates
\[
 (x,z)=\bigl(K_{s,r}^{-1}w, V_r^*K_{s,r}^{-1}w\bigr),
\qquad
 \|P_{\min}\|^2=w^*K_{s,r}^{-1}w.
 \tag{IFB10}
\]
Indeed $x+V_rz=K_{s,r}K_{s,r}^{-1}w=w$. The displayed vector
is orthogonal to every $(-V_rz',z')$, since its inner product with
that vector is zero. Every other lift differs by such a relation
vector, so the Pythagorean identity proves uniqueness and the
minimum formula. It follows in the original fixed remainder frame
that
\[
 G_{q+r-1}^{(s)}=C_s^{-*}K_{s,r}^{-1}C_s^{-1},\qquad
 \det G_{q+r-1}^{(s)}
       =\frac{\prod_{n=0}^{q-1}h_n^{(s)}}{\Omega_{s,r}}.
 \tag{IFB11}
\]
This gives the entire quotient Gram, including every off-diagonal
entry, before its determinant is taken.

\subsection{The relation determinant as a positive finite sum}

The integer $q$ is divisible by four. Pairing both original signed
root coordinates in IFB1 gives the exact real polynomial
\[
 \vartheta(y):=\chi(c+iy)
 =\prod_{A\in\{\delta,3\delta,\ldots,k\delta\}}
   \prod_{B\in\{\gamma,3\gamma,\ldots,k\gamma\}}
 \bigl[((y-B)^2+A^2)((y+B)^2+A^2)\bigr]^e.
 \tag{IFB12}
\]
Each quartet contributes two negative factors from its two opposite
real-root pairs, so their product has the positive sign displayed.
Each factor is strictly positive on $\mathbb R$ because $A>0$.
Thus $\vartheta$ is real, even, strictly positive on the real line,
and monic of degree $q$; its squared modulus is exactly
$W_\chi(y)=\vartheta(y)^2$. Write
\[
 \vartheta(y)=\sum_{a=0}^q c_a y^a,\qquad
 W_\chi(y)=\sum_{d=0}^{2q}w_dy^d,\qquad
 w_d=\sum_{a+a'=d}c_ac_{a'}.
 \tag{IFB13}
\]
These coefficients retain all roots and all repetitions. For an
explicit finite rule producing them, each factor of IFB12 is
$(y^4+2(A^2-B^2)y^2+(A^2+B^2)^2)^e$. Expand it as
\[
 \sum_{u+v+w=e}
 \frac{e!}{u!v!w!}\,[2(A^2-B^2)]^v(A^2+B^2)^{2w}
             y^{4u+2v},
 \tag{IFB14}
\]
and multiply these finite coefficient lists. In particular this rule
does not discard their signed middle coefficients.

Let $T_{n,h}$ be the coefficient of $p_h$ in $y^n$. The exact
finite recurrence following from IFB5 is
\[
 T_{0,0}=1,\quad T_{0,h}=0\ (h\ne0),\quad
 T_{n+1,h}=T_{n,h-1}+(h+1)(b_s+h)T_{n,h+1},
 \tag{IFB15}
\]
with indices outside $0\leq h\leq n$ interpreted as zero. Induction
using $yp_h=p_{h+1}+h(b_s+h-1)p_{h-1}$ proves the formula and
$y^n=\sum_hT_{n,h}p_h$. In particular $T_{n,n}=1$, and only
$h\equiv n\pmod2$ can occur. Define the full rectangular matrix
\[
 A_{h,j}=\sum_{a=0}^q c_aT_{a+j,h},\quad
 0\leq h\leq q+r-1,\quad 0\leq j<r,
 \qquad
 L_{h,j}=\sqrt{h_h^{(s)}}A_{h,j}.
 \tag{IFB16}
\]
Here $h_h^{(s)}$ means the source norm with index $h$; the matrix
letter $L$ denotes the full weighted coefficient matrix only in this
section. The definition says exactly that
$\vartheta(y)y^j=\sum_h A_{h,j}p_h(y)$.
Thus the Gram of these $r$ polynomials is $L^*L$.

In the original $S$-coefficient frame the relation vectors are
$(\chi,\chi S,\ldots,\chi S^{r-1})$. Its coefficient change to
the vectors $(\vartheta,\vartheta y,\ldots,\vartheta y^{r-1})$
is the exact triangular matrix
\[
 (T_c)_{h,j}=\binom jh c^{j-h}i^h\quad(0\leq h\leq j<r).
 \tag{IFB17}
\]
Its inverse is substitution $y=(S-c)/i$, and
$|\det T_c|=1$. Consequently the original relation determinant
$Z_{s,r}:=\det\operatorname{Gram}_{m_{0,s}}
(\chi,\ldots,\chi S^{r-1})$ equals $\det(L^*L)$.
We set $Z_{s,0}=1$.

Let $L_{\mathrm{low}}$ comprise rows $0,\ldots,q-1$ of $L$ and
$R$ comprise rows $q,\ldots,q+r-1$. Since $\vartheta y^j$ is
monic of degree $q+j$, the matrix $R$ is upper triangular with
diagonal $\sqrt{h_{q+j}^{(s)}}$. Hence it is invertible and
\[
 |\det R|^2=\prod_{j=0}^{r-1}h_{q+j}^{(s)},\qquad
 L_{\mathrm{low}}R^{-1}=-V_r.
 \tag{IFB18}
\]
For the second identity, the column space of $L$ is the relation
space in IFB9. Its high coordinates are precisely $R$ times its
coefficient vector; expressing a relation by its high coordinate
therefore gives the low graph matrix $L_{\mathrm{low}}R^{-1}$.
Uniqueness of the graph in IFB9 proves the asserted sign and every
matrix entry. Therefore
\[
 Z_{s,r}
 =\left(\prod_{j=0}^{r-1}h_{q+j}^{(s)}\right)
       \det(I_r+V_r^*V_r)
 =\left(\prod_{j=0}^{r-1}h_{q+j}^{(s)}\right)\Omega_{s,r}.
 \tag{IFB19}
\]
To justify the second equality without discarding zero singular
directions, eliminate either diagonal identity block in
$\left(\begin{smallmatrix}I_q&V_r\\-V_r^*&I_r\end{smallmatrix}\right)$.
The two block eliminations have determinant one and give
$\det(I_q+V_rV_r^*)=\det(I_r+V_r^*V_r)$, for all dimensions.

The full finite expansion of the same determinant is
\[
 \boxed{\displaystyle
 \Omega_{s,r}=
 \sum_{\substack{I\subset\{0,\ldots,q+r-1\}\\|I|=r}}
 |\det A_{I,\{0,\ldots,r-1\}}|^2
 \frac{\prod_{h\in I}h_h^{(s)}}
      {\prod_{j=0}^{r-1}h_{q+j}^{(s)}}.}
 \tag{IFB20}
\]
For completeness, expand each column of $L$ in the standard
orthonormal coordinates of $\mathbb C^{q+r}$. Their exterior
product has, in the exterior frame indexed by $r$-element subsets
$I$, coefficient $\det L_I$. The exterior frame is orthonormal,
so its squared norm is $\sum_I|\det L_I|^2$. Direct expansion of
that same squared norm gives $\det(L^*L)$ by the permutation
formula for a determinant. Substituting IFB16 and dividing by the
explicit positive factor in IFB19 proves IFB20. In particular the
subset $I=\{q,\ldots,q+r-1\}$ contributes exactly one. All other
terms are nonnegative; none of their signed coefficient calculations
has been replaced by a degree estimate. When $r=0$, the single
empty minor contributes one.

There is also a direct moment evaluation of exactly the same mass.
Define polynomials $P_n(z)$ by
\[
 P_0(z)=1,\qquad
 P_{n+1}(z)=b_szP_n(z)+(1+z^2)P_n'(z).
\]
Induction by differentiation gives
$\frac{d^n}{dt^n}(\cos t)^{-b_s}
=(\cos t)^{-b_s}P_n(\tan t)$. Differentiating IFB3 at zero
is justified by its exponential bound. It follows that
\[
 \boxed{\displaystyle
 Z_{s,r}=M_s^r\det\left[
       \sum_{d=0}^{2q}w_dP_{d+i+j}(0)
                          \right]_{i,j=0}^{r-1}.}
 \tag{IFB21}
\]
The moment degree required at the largest original endpoint $r=q+1$
is exactly $2q+2r-2=4q$. The finite recurrences IFB14--16 and
IFB21 therefore calculate every contribution from the original
parameters in finitely many additions, multiplications and divisions
by the displayed nonzero positive constants. IFB20 certifies the
positivity of the result even when coefficients in IFB21 cancel.

\subsection{The actual spectral-sum mass at the four ranks}

For $r\geq1$, let $\Delta(y)=\prod_{1\leq i<j\leq r}(y_j-y_i)$
and define the actual exterior measure and actual sum map
\[
 d\mathfrak m_{s,\chi,r}(y_1,\ldots,y_r)
 =\frac1{r!}\Delta(y)^2\prod_{j=1}^r
       W_\chi(y_j)\,dm_{0,s}(y_j),\qquad
 \Sigma_r(y_1,\ldots,y_r)=\sum_{j=1}^ry_j.
 \tag{IFB22}
\]
Set $\nu_{s,\chi,r}=(\Sigma_r)_*\mathfrak m_{s,\chi,r}$.
At $r=0$ the source space is a one-point space of mass one,
$\Sigma_0=0$, and $\nu_{s,\chi,0}=\delta_0$.
Expanding both Vandermonde determinants in IFB22 gives
\[
 \nu_{s,\chi,r}(\mathbb R)=Z_{s,r}.
 \tag{IFB23}
\]
To check the prefactor, write $\Delta=\det[y_i^{j-1}]$ and
expand it twice. Integrating a product gives a product of the
moments of $W_\chi dm_{0,s}$. Holding the first permutation fixed
and reindexing the second gives the determinant of its $r$ by $r$
moment matrix; the first permutation has $r!$ possible values.
Those $r!$ values cancel precisely the $1/r!$ in IFB22. Absolute
integrability follows from the polynomial degree and the exponential
source bound. The moment determinant agrees with the original
$S$-relation determinant by IFB17. This proves IFB23 with its
entire original mass.

Here is the precise retained fibre map. Write $\rho_s$ for the
density in IFB2 and, for $r\geq2$, put
\[
 y_r=Y-\sum_{a=1}^{r-1}y_a,\qquad
 F_r(Y;y_1,\ldots,y_{r-1})
 =\frac1{r!}\Delta(y)^2\prod_{a=1}^r W_\chi(y_a)\rho_s(y_a),
\]
\[
 g_r(Y)=\int_{\mathbb R^{r-1}}F_r(Y;y')\,dy'.
 \tag{IFB24}
\]
The change of variables $(y_1,\ldots,y_r)\mapsto(y_1,\ldots,
y_{r-1},Y)$ has Jacobian determinant one. Therefore
$d\nu_{s,\chi,r}(Y)=g_r(Y)dY$. These integrals are finite for every
$Y$: a polynomial in $|Y|+\sum|y_a|$ times
$\exp(-\theta\sum_{a<r}|y_a|)$ is an integrable dominating
function when $Y$ is fixed. Each $g_r(Y)>0$, because the source
densities and $W_\chi$ are positive and the distinct-coordinate
configurations form a nonempty open subset of the fibre.
For $r=1$, the formula is $g_1(Y)=W_\chi(Y)\rho_s(Y)>0$.

The map
\[
 \mathcal U_r:L^2(\nu_{s,\chi,r})\longrightarrow
                 L^2(\mathfrak m_{s,\chi,r}),\qquad
 (\mathcal U_r f)(y)=f(\Sigma_r y)
 \tag{IFB25}
\]
is an isometry, directly by the defining integral of a pushforward.
Its full adjoint, for $r\geq2$, is
\[
 (\mathcal U_r^*H)(Y)
 =\frac1{g_r(Y)}\int H(y_1,\ldots,y_{r-1},Y-\textstyle\sum y_a)
                         F_r(Y;y')\,dy'.
 \tag{IFB26}
\]
The integral exists for almost every $Y$ by Cauchy--Schwarz in the
fibre and Fubini. The same inequality shows its squared modulus,
multiplied by $g_r(Y)$, is bounded by the fibre integral of
$|H|^2F_r$; hence the adjoint is a contraction. Multiplying by a
test function $f(Y)$ and applying Fubini proves the adjoint identity
and thus IFB26. Its kernel is exactly the functions with weighted
fibre integral zero for almost every $Y$. For $r=1$, both maps are
the identity; at $r=0$ they are the identity of the one-dimensional
one-point $L^2$ space. Thus the fibre kernel and the full source
norm have specified domains and codomains; passing to a sum does
not silently identify the entire exterior space with its image.

The original complex spectral sum is
$\sum_{a=1}^r(c+iy_a)=rc+iY$. Accordingly the actual target line
of this pushforward is $rc+i\mathbb R$. If a Gamma density of
order $a$ is written on its natural line $a/2+i\mathbb R$, the
coordinate isomorphism for that constituent is
\[
 a/2+iY\longmapsto rc+iY,
 \quad\hbox{translation by }rc-a/2.
 \tag{IFB27}
\]
It preserves $dY$ and the complete density and mass. Different
constituents use their own stated translation; none changes the
original target centre $rc$.

Let $\mathcal D_{s,\chi,r}(z)$ denote the polynomial provided by
the intrinsic exterior/sum calculation, so its exact Laplace
transform is
\[
 \int e^{tY}\,d\nu_{s,\chi,r}(Y)
 =M_s^r(\cos t)^{-[b_sr+r(r-1)]}
             \mathcal D_{s,\chi,r}(\tan t).
 \tag{IFB28}
\]
This polynomial is completely determined by IFB21 as a differentiated
moment determinant, or by the full Wronskian construction in the
intrinsic source provider. To check this definition directly, let
$Q_0(z)=\sum_dw_dP_d(z)$ and
$Q_{j+1}(z)=b_szQ_j(z)+(1+z^2)Q_j'(z)$.
The Laplace version of IFB23 has moment entries
$M_s(\cos t)^{-b_s}Q_{i+j}(\tan t)$.
Thus one may take
\[
 \mathcal D_{s,\chi,r}(z)
 =(1+z^2)^{-r(r-1)/2}\det[Q_{i+j}(z)]_{i,j<r}.
\]
This expression is a polynomial: its determinant is the
$t$-Wronskian of the first $r$ derivatives of
$(\cos t)^{-b_s}Q_0(\tan t)$, and the following elementary two
identities show the factors explicitly.
For a common factor $f$, successive row subtractions in
$[\partial_t^i(fg_j)]$ give determinant $f^r$
times the Wronskian of the $g_j$; the subtraction coefficients come
from the product rule, and the triangular derivative matrix has
diagonal $f$. For a coordinate $z=z(t)$, the chain rule expresses
$\partial_t^i$ as $(z')^i\partial_z^i$ plus lower derivatives;
the determinant multiplier is $(z')^{r(r-1)/2}$.
Apply these identities with $f=(\cos t)^{-b_s}$,
$g_j=Q_j(z)$ and $z'=1+z^2$ to obtain
$\mathcal D_{s,\chi,r}(z)=\det[\partial_z^iQ_j(z)]_{i,j<r}$.
This also proves IFB28 without any assumed statement about the
degree or signs of this polynomial. Its full degree and Gamma
constituents are evaluated in the intrinsic source provider.

At the original four ranks its constant term is now calculated by
the positive original relation graph, not by size counting:
\[
 \boxed{\displaystyle
 \mathcal D_{s,\chi,r}(0)
 =\left(\prod_{j=0}^{r-1}(q+j)!(b_s)_{q+j}\right)
                       \Omega_{s,r},\quad r=0,1,q,q+1.}
 \tag{IFB29}
\]
Indeed IFB28 at $t=0$ gives $M_s^r\mathcal D(0)=Z_{s,r}$.
Substituting IFB19 and every factor
$h_{q+j}^{(s)}=M_s(q+j)!(b_s)_{q+j}$ proves IFB29. In particular
the full exterior-source mass appears to power $r$ on both sides
before its algebraic cancellation. Equation IFB29 is the exact
factor relating the intrinsic Gamma expression to the original
quotient determinant.

\subsection{The original four-endpoint identity and its positive update chain}

Write $\mathcal B_k^{(1)}=\mathcal B_k^\sigma$ and
$\mathcal B_k^{(k)}=\mathcal B_k^0$. The monic source norms in the
original $S$-basis are $h_n^{(s)}$: multiplication of $u_n$ by
$i^n\sqrt{h_n^{(s)}}$ makes it monic and leaves its squared norm
$h_n^{(s)}$. The triangular monic source/relation frame in BSL2--3
therefore gives, at the original four degrees, the full signed formula
\[
 \begin{split}
 \mathcal B_k^{(s)}
 &=S_q^{(s)}+\log Z_{s,q}+\log Z_{s,q+1}-\log Z_{s,1},\\
 S_q^{(s)}
 &=-\log h_q^{(s)}-2\sum_{j=q+1}^{2q-1}\log h_j^{(s)}
                                      -\log h_{2q}^{(s)}.
 \end{split}
 \tag{IFB30}
\]
The omitted rank-zero logarithm is exactly $\log Z_{s,0}=0$.
One can also prove this directly from IFB11 by taking the prescribed
$+,+,-,-$ signs at $r=0,1,q,q+1$. Substituting IFB19 into the
displayed formula cancels the monic norm factors as follows:
the two high products contain $h_q^2$, each
$h_{q+1},\ldots,h_{2q-1}$ twice, and $h_{2q}$ once; division by
the low factor $h_q$ leaves precisely the factors canceled by
$S_q^{(s)}$. In particular the full source mass exponent in the
relation terms is $q+(q+1)-1=2q$, while that in $S_q^{(s)}$ is
$-1-2(q-1)-1=-2q$. The exact result is
\[
 \boxed{\mathcal B_k^{(s)}
    =\log\Omega_{s,q}+\log\Omega_{s,q+1}-\log\Omega_{s,1}.}
 \tag{IFB31}
\]
Every original sign and all low and high polynomial factors have
now been retained and evaluated before this cancellation.

Define the finite, strictly positive update factors
\[
 d_{s,j}=1+v_j^*K_{s,j}^{-1}v_j,\qquad j\geq0.
 \tag{IFB32}
\]
Since $K_{s,j}=I_q+\sum_{a<j}v_av_a^*$ is positive definite, its
inverse exists and $d_{s,j}\geq1$. Multiplying
$K_{s,j+1}=K_{s,j}+v_jv_j^*$ by $K_{s,j}^{-1/2}$ on both sides
reduces its determinant ratio to that of $I+ww^*$,
$w=K_{s,j}^{-1/2}v_j$. In an orthonormal frame with first vector
$w/\|w\|$ when $w\ne0$, that operator has eigenvalues
$1+\|w\|^2,1,\ldots,1$; the case $w=0$ is the identity.
Consequently
\[
 \Omega_{s,r}=\prod_{j=0}^{r-1}d_{s,j},\qquad
 K_{s,j+1}^{-1}=K_{s,j}^{-1}
 -\frac{K_{s,j}^{-1}v_jv_j^*K_{s,j}^{-1}}{d_{s,j}}.
 \tag{IFB33}
\]
The inverse identity follows by multiplication by
$K_{s,j}+v_jv_j^*$; its two rank-one terms cancel exactly.
Starting with $K_{s,0}^{-1}=I_q$, it is a finite formula for every
inverse required in IFB32. The original four-endpoint expression
therefore equals the complete nonnegative sum
\[
 \boxed{\displaystyle
 \mathcal B_k^{(s)}
 =\log d_{s,0}+2\sum_{j=1}^{q-1}\log d_{s,j}+\log d_{s,q}.}
 \tag{IFB34}
\]
This is proved from the signed identity IFB31: each factor with
$1\leq j<q$ occurs in both high endpoints, $j=q$ occurs once,
and the rank-one low endpoint removes one of the two occurrences
of $j=0$. Formula IFB34 preserves all singular directions and all
off-diagonal interactions through the full inverse $K_{s,j}^{-1}$.
It is not an assertion that these interactions are independent.

\subsection{A finite companion calculation of every update vector}

Let $E_{\vartheta}=\mathbb C[y]/(\vartheta)$ in its monomial
frame and let $Y_\vartheta$ be the matrix of multiplication by $y$.
Thus its column $j<q-1$ is the vector with a one in row $j+1$,
and its last column is $(-c_0,\ldots,-c_{q-1})^t$.
The exact algebra isomorphism is
\[
 E_{\vartheta}\longrightarrow E_\chi,\qquad
 [P(y)]\longmapsto[P((S-c)/i)],
 \tag{IFB35}
\]
because substituting $y=(S-c)/i$ in
$\vartheta(y)=\chi(c+iy)$ gives $\chi(S)$, and the inverse
substitution is $S=c+iy$. Let $E_s$ be the coefficient matrix of
$p_0/\sqrt{h_0^{(s)}},\ldots,p_{q-1}/\sqrt{h_{q-1}^{(s)}}$
in the $y$-monomial frame; its diagonal is
$1/\sqrt{h_n^{(s)}}$. Hence $E_s$ is invertible. Define
\[
 Y_s=E_s^{-1}Y_\vartheta E_s,\qquad
 a_n=\sqrt{n(b_s+n-1)}\ (n\geq1),\quad a_0=0.
 \tag{IFB36}
\]
All entries are now determined by the finite original coefficients
in IFB14 and the finite original recurrence IFB5. For all $n\geq0$
let $z_n=C_s^{-1}[u_n]$. Its initial values are
$z_{-1}=0$ and $z_n$ the $n$th coordinate vector for $0\leq n<q$.
Dividing IFB5 by $\sqrt{h_n^{(s)}}$ and using
$h_n^{(s)}/h_{n-1}^{(s)}=n(b_s+n-1)$ gives
\[
 \boxed{\displaystyle
 z_{n+1}=\frac{Y_sz_n-a_nz_{n-1}}{a_{n+1}},\qquad
 v_j=z_{q+j}.}
 \tag{IFB37}
\]
Multiplication in the quotient is exactly $Y_s$ by IFB35--36, so
this is a recurrence for the actual remainders, including all of
their original polynomial relation coefficients. Every denominator
$a_{n+1}$ is positive. It suffices to run it through $n=2q-1$ to
obtain the last vector $v_q=z_{2q}$ used in IFB34. Combining
IFB14, IFB36--37, and IFB32--33 is thus a finite exact calculation
of the original four-endpoint quantity, using a fixed $q$-dimensional
space and its full matrices.

As an explicit coefficient check, summing the coefficients of
$y^{q-2}$ in the original quartets gives
\[
 c_{q-2}=\frac{qk(k+2)(\delta^2-\gamma^2)}6.
 \tag{IFB38}
\]
Indeed $\sum_{a\in\{1,3,\ldots,k\}}a^2=k(k+1)(k+2)/6$ and
there are $(k+1)/2$ choices of the other positive coordinate.
Each quartet contributes $2e(A^2-B^2)$, which proves IFB38.
The coefficient of $y^{n-2}$ in $p_n$ is
\[
 -A_n^{(s)},\qquad
 A_n^{(s)}=\sum_{j=1}^{n-1}j(b_s+j-1)
          =\frac{n(n-1)(3b_s+2n-4)}6.
 \tag{IFB39}
\]
The first equality follows successively from the coefficient of
$y^{n-1}$ in the recurrence for $p_{n+1}$, starting with $p_1=y$;
the second uses the finite sums of $j$ and $j^2$. Consequently
the coefficient of $p_{q-2}$ in the full original $\vartheta$
is exactly $c_{q-2}+A_q^{(s)}$. In IFB20 the minor obtained by
replacing the row $q$ in the highest-degree minor by $q-2$ is
therefore $c_{q-2}+A_q^{(s)}$, for every $r\geq1$.
To check that assertion, its first column has this entry in the
replacement row and is zero in all remaining rows, while the
remaining top-degree block is upper triangular with diagonal one.
It follows that
\[
 \Omega_{s,r}\geq
 1+\frac{(c_{q-2}+A_q^{(s)})^2h_{q-2}^{(s)}}{h_q^{(s)}}
 \quad(r\geq1).
 \tag{IFB40}
\]
For $r\geq2$, the analogous replacement of row $q+1$ by $q-1$
has squared minor $(c_{q-2}+A_{q+1}^{(s)})^2$: its two lowest
rows and first two columns have determinant
$-(c_{q-2}+A_{q+1}^{(s)})$, and the remaining block is upper
triangular. Making both replacements gives the product of the
two coefficients, because their opposite-parity cross entries
vanish. Keeping these four distinct nonnegative summands in
IFB20 proves the stronger finite bound
\[
 \Omega_{s,r}\geq
 \left(1+\frac{(c_{q-2}+A_q^{(s)})^2h_{q-2}^{(s)}}{h_q^{(s)}}\right)
 \left(1+\frac{(c_{q-2}+A_{q+1}^{(s)})^2h_{q-1}^{(s)}}
                         {h_{q+1}^{(s)}}\right),\quad r\geq2.
 \tag{IFB41}
\]
These are specific calculated terms in the full original polynomial
mass; IFB20 and IFB34 continue to retain every remaining term.
In particular they assert no asymptotic replacement of the full
quantity by these selected summands.

Finally $\vartheta$ is even and each $p_n$ has its original parity.
Division by this even monic polynomial preserves parity. Hence
$v_j$ has coordinates only in indices congruent to $q+j$, and
$K_{s,r}$ has zero entries between the two source parity subspaces.
This follows entrywise from IFB8, not by removing coordinates.
The full matrix $C_s$ in IFB11 carries this exact structure back
to the original $S$-remainder frame. The downstream GMB calculation
then composes that same map with the original mixed boundary map;
its contractions use the entire matrix $K_{s,r}$ and all its
original source phases.
